% ─────────────────────────────────────────────────────────────────────────────
% CHAPTER 2: Literature Survey
% ─────────────────────────────────────────────────────────────────────────────

This chapter reviews relevant academic literature, commercial systems, and existing technical approaches related to online grocery platforms, recommendation systems, demand forecasting, autonomous agents, and reverse auction marketplaces.

\section{Evolution of Online Grocery and Quick-Commerce}

In recent years, the online grocery market has shifted from scheduled next-day deliveries (such as traditional BigBasket or Amazon Fresh) toward hyper-local quick-commerce platforms like Blinkit, Zepto, and Swiggy Instamart. These platforms rely on geographically distributed micro-warehouses called ``dark stores,'' each holding between two thousand and four thousand common stock-keeping units (SKUs) located within a three-kilometer radius of urban neighborhoods.

While this model achieves very low delivery times, studies on retail logistics note several structural constraints:
\begin{itemize}
    \item \textbf{Inventory Limitations:} A single dark store cannot maintain the variety of items found across an entire city or district. When niche brands, bulk bags, or seasonal local produce are needed, the customer faces out-of-stock messages.
    \item \textbf{Pricing Inflexibility:} Prices in centralized dark stores are determined by algorithmic markup rules with no room for buyer negotiation or local merchant competition.
    \item \textbf{Exclusion of Local Neighborhood Shops:} Independent local merchants, who often possess superior connections with regional farmers and wholesale grain markets, are largely excluded from modern quick-commerce infrastructure.
\end{itemize}

\section{Recommendation Systems in Retail}

Personalized product recommendation is widely recognized as a primary driver of customer retention in e-commerce.

\subsection{Collaborative Filtering}
Collaborative filtering (CF), introduced formally by Resnick et al.~\cite{resnick1994grouplens} and expanded by Sarwar et al.~\cite{sarwar2001item}, predicts user preferences by identifying patterns among similar users or items. In item-based CF, the similarity between two items $i$ and $j$ is computed across the user rating vector using the cosine similarity metric:
\begin{equation}
    \text{sim}(i, j) = \frac{\sum_{u \in U} R_{u,i} \cdot R_{u,j}}{\sqrt{\sum_{u \in U} R_{u,i}^2} \cdot \sqrt{\sum_{u \in U} R_{u,j}^2}}
\end{equation}
While collaborative filtering works well when substantial interaction data is available, it suffers from the well-documented \textbf{cold-start problem}: when a new product is added or a new customer signs up, there are no prior interactions, causing collaborative filtering alone to fail.

\subsection{Content-Based Filtering}
Content-based filtering, discussed by Pazzani and Billsus~\cite{pazzani2007content}, avoids the cold-start problem by analyzing product features directly (such as category, brand, dietary attributes, and price tier). A user profile vector is constructed from the attributes of items the user has previously interacted with, and candidate products are scored by vector similarity. However, content-based methods tend to suffer from \textbf{overspecialization}, repeatedly recommending items very similar to what the user already purchased.

\subsection{Hybrid Approaches}
Burke~\cite{burke2002hybrid} showed that combining collaborative and content-based approaches in an ensemble overcomes the individual weaknesses of both methods. In GroDealz, we implement a weighted hybrid ensemble that also incorporates contextual signals, such as the time of day (e.g., boosting milk and bakery products in the morning, and dinner vegetables and rice in the evening).

\section{Demand Forecasting in Perishable Retail}

Accurate demand prediction is essential for grocery stores because fresh items like vegetables, fruits, and dairy have very short shelf lives. Overstocking causes food spoilage and financial loss, while understocking leads to lost sales and disappointed customers.

Classic time-series methods such as the Box-Jenkins ARIMA models~\cite{box2015time} require stationary data and significant computational overhead. For hyper-local stores operating on standard web servers, simpler lightweight techniques are often more practical:
\begin{itemize}
    \item \textbf{Exponential Smoothing:} Gives exponentially decreasing weights to older observations, allowing recent order bursts to adjust the forecast quickly.
    \item \textbf{Moving Average Trend:} Tracks baseline demand over a multi-day rolling window to capture underlying weekly cycles.
    \item \textbf{Multivariate Linear Regression:} Relates order volume directly to calendar features such as hour of the day, day of the week, and weekend indicators.
\end{itemize}
Armstrong~\cite{armstrong2001combining} established that combining separate forecasting models into an ensemble frequently outperforms individual models. GroDealz adopts this insight by averaging exponential smoothing with linear regression to generate hourly per-product demand forecasts.

\section{Autonomous Agents and Multi-Step Task Execution}

Recent breakthroughs in Large Language Models (LLMs) and agentic frameworks have fundamentally altered automated decision systems in e-commerce. Early conversational systems relied either on scripted decision trees (Weizenbaum~\cite{weizenbaum1966eliza}) or intent-classification chatbots with rigid slot-filling (F{\o}lstad and Brandtz{\ae}g~\cite{folstad2017chatbots}). While effective for narrow FAQs, such systems cannot perform multi-step planning, stateful backtracking, or dynamic tool execution.

The introduction of the ReAct paradigm by Yao et al.~\cite{yao2022react} demonstrated that prompting language models to alternate between explicit reasoning traces (``Thought'') and executable environment actions (``Act'') significantly improves performance on complex multi-hop queries. In hyper-local grocery shopping, an agent cannot simply answer a question; it must execute an autonomous operational loop:
\begin{enumerate}
    \item \textbf{Structured Parameter Extraction:} Parse informal colloquial natural language or Hinglish into validated domain entities (family size $N$, duration $D$, target budget $B$, dietary tags).
    \item \textbf{Parallel Environment Discovery:} Asynchronously query multiple distributed merchant databases to retrieve live product candidates, prices, and stock indicators.
    \item \textbf{Data Normalization and Constraint Verification:} Standardize disparate unit measurements (e.g., converting 500~g packets and 1~kg sacks into unified unit-cost metrics) and enforce hard budget constraints.
    \item \textbf{Combinatorial Basket Optimization:} Formulate and solve the single-store versus multi-store optimization tradeoff, balancing per-item price savings against individual store delivery charges.
    \item \textbf{Auditable Itinerary Generation:} Compile an inspectable, human-verifiable store picking audit that allows the buyer to review individual store contributions before order placement.
\end{enumerate}

Traditional agent frameworks like LangChain standard chains execute purely linear directed acyclic graphs (DAGs), which fail when an agent encounters unexpected real-world edge cases (such as a chosen item being out of stock across all preferred stores). To support iterative refinement and conditional recovery, GroDealz employs \textbf{LangGraph}, a cyclic state-machine framework where agent execution is modeled as a state graph with typed state channels, conditional branching, and deterministic fallback nodes. By constraining language model outputs to strictly typed Pydantic data schemas, the system eliminates mathematical hallucination, ensuring that all item quantities and financial figures reflect verified database records.

\section{Reverse Auctions and Multi-Attribute Request for Quotation (RFQ) Systems}

In standard forward e-commerce marketplaces (such as Amazon or Blinkit), price discovery is strictly one-directional: the vendor sets a take-it-or-leave-it price, and the customer has no channel to express willingness-to-pay or request customized volume discounts. In contrast, reverse auction mechanisms invert this market dynamic: the buyer announces a specific requirement, and multiple pre-qualified suppliers submit competitive bids to fulfill the order.

Extensive literature on auction theory (De Vries and Vohra~\cite{de2003combinatorial}) distinguishes between several reverse bidding formats:
\begin{itemize}
    \item \textbf{First-Price Sealed-Bid Reverse Auctions:} Suppliers submit independent, private bids without visibility into competing offers. The buyer awards the contract to the lowest bidder. While computationally simple, sealed-bid formats often induce the ``winner's curse'' and prevent suppliers from dynamically adjusting prices based on real-time capacity.
    \item \textbf{Iterative English Reverse Auctions:} Suppliers openly observe competing bids and progressively lower their prices in descending rounds. While driving prices down to marginal costs, open descending auctions can trigger predatory bidding wars and discourage small merchant participation.
    \item \textbf{Multi-Attribute Interactive Quotation Negotiation:} In consumer and institutional procurement, price is rarely the sole evaluation criterion. Delivery reliability, fulfillment lead time, merchant reputation, and itemized product brand quality are equally vital. Multi-Attribute Utility Theory (MAUT) provides a mathematical basis for evaluating multi-dimensional bids by constructing a composite scoring function:
    \begin{equation}
        U(\vec{x}) = \sum_{k=1}^{K} w_k \cdot v_k(x_k)
    \end{equation}
    where $w_k$ denotes the normalized importance weight of attribute $k$, and $v_k(x_k)$ is the attribute-specific value function.
\end{itemize}

GroDealz implements an itemized multi-attribute quotation engine (Module B). Rather than bidding arbitrary lump sums, merchants submit granular line-item prices alongside explicit delivery timestamps and stock commitments. Furthermore, the platform empowers the customer to choose their immediate evaluation priority dynamically---prioritizing lowest cost, fastest delivery speed, highest vendor reputation, or a balanced 100-point composite score. This guarantees transparent price formation while preventing merchants from concealing inflated delivery or packaging fees inside opaque bundles.

\section{Data Security, Privacy, and Transactional Concurrency in Multi-Sided Platforms}

Multi-sided digital marketplaces handle concurrent interactions between competing merchants and buyers, introducing critical data protection and concurrency challenges:

\begin{enumerate}
    \item \textbf{Personally Identifiable Information (PII) Protection in Customer Support:} When buyers submit support tickets or interact with automated AI CRM systems regarding delivery disputes, they frequently share sensitive details such as contact telephone numbers, home delivery landmarks, and transaction identifiers. To comply with modern data privacy principles, sensitive fields must be sanitized before being transmitted to external language model APIs or centralized analytics logs. Implementing rule-based regex masking on the client and server boundary ensures that phone numbers and email addresses are replaced with anonymized tokens (e.g., \texttt{[PHONE\_REDACTED]}) while preserving the semantic context required for intent classification.
    \item \textbf{Concurrency and Atomic State Transitions in Reverse RFQs:} In an active reverse marketplace, multiple events occur asynchronously: merchants submit bids, buyers issue counter-offers, and quotations approach expiration timestamps. A critical concurrency hazard arises when a customer accepts a quotation: the selected quote must transition to \texttt{accepted}, all other active quotations for that request must be atomically updated to \texttt{rejected}, a confirmed \texttt{Order} document must be instantiated, and the parent \texttt{BuyerRequest} status must become \texttt{converted}. If these operations are executed non-atomically, race conditions could allow a merchant to revise a quote simultaneously or enable double-acceptance. Enforcing atomic database document updates and transactional validation prevents state inconsistency.
    \item \textbf{Role-Based Cryptographic Access Control:} Because merchants and customers share the same underlying API endpoints, the application server must cryptographically verify user identity and authority on every request. JSON Web Tokens (JWT) signed with high-entropy HMAC-SHA256 secrets provide stateless, tamper-proof session tokens containing encrypted user identifiers and operational roles, preventing unauthorized access across disparate merchant dashboards.
\end{enumerate}

\section{Comparative Summary}

Table~\ref{tab:lit_comparison} highlights how GroDealz compares with standard existing platforms across essential functional areas.

\begin{table}[H]
\centering
\caption{Comparison of GroDealz with Existing Grocery Delivery Approaches}
\label{tab:lit_comparison}
\begin{tabular}{>{\raggedright\arraybackslash}p{3.2cm} >{\raggedright\arraybackslash}p{2.8cm} >{\raggedright\arraybackslash}p{3.2cm} >{\raggedright\arraybackslash}p{4.2cm}}
\toprule
\textbf{Feature} & \textbf{Quick Commerce} & \textbf{Traditional Grocery} & \textbf{GroDealz} \\
\midrule
Store Model & Single Dark Store & Central Warehouse & Multi-Store Local Network \\
\addlinespace
Autonomous AI Planner & None (Search Only) & None (Search Only) & Natural Language Multi-Store Agent \\
\addlinespace
Basket Optimization & No Comparison & Single Basket & Multi-Store vs Single-Store Net Savings \\
\addlinespace
Reverse RFQ Marketplace & Not Supported & Not Supported & Full Itemized Merchant Bidding (Module B) \\
\addlinespace
Bid Ranking Modes & Fixed Price & Fixed Price & Lowest Price, Fastest Delivery, Rating, Balanced \\
\addlinespace
Live Store Journey Audit & None & None & Visual Picking Audit and Savings Proof \\
\addlinespace
Real-Time Tracking & Polling / SMS & Delivery Window & Live WebSockets (Socket.IO) \\
\bottomrule
\end{tabular}
\end{table}
